\chapter{Ruang Banach dan Teorema Fundamentalnya}
\label{ch:b3:banach}

Analisis fungsional mempelajari ruang bernorma yang berdimensi tak berhingga
lewat operator dan fungsional yang hidup di atasnya. Penemuan
pendirinya adalah bahwa \emph{\omterm{def:b3:complete:complete}{kelengkapan}}, lewat teorema
Baire, memaksa keseragaman yang kuat: keluarga operator yang terbatas titik demi
titik ternyata terbatas normanya (Banach--Steinhaus),
bijeksi \omterm{def:b3:topology:continuity}{kontinu} berinvers \omterm{def:b3:topology:continuity}{kontinu} (pemetaan terbuka), dan grafik
mendeteksi \omterm{def:b3:topology:continuity}{kekontinuan} (grafik tertutup). Pilar yang lain,
\emph{Hahn--Banach}, tak memerlukan \omterm{def:b3:complete:complete}{kelengkapan} sama sekali --- hanya
lema Zorn --- dan menjamin bahwa \omterm{def:b3:banach:operator}{ruang dualnya} cukup kaya
untuk melihat setiap vektor. Bab ini membuktikan keempat teorema itu dan
mengujinya pada ruang barisan klasik $\ell^p$, pada
perhitungan \omterm{def:b3:banach:operator}{dual} yang konkret, dan pada terapan yang sungguh
mengejutkan: bahwa ada fungsi \omterm{def:b3:topology:continuity}{kontinu} berperiode $2\pi$
yang deret Fouriernya \emph{divergen} di sebuah titik --- sehingga menjawab
secara negatif pertanyaan yang ditinggalkan terbuka pada Tahun ke-2.

Sepanjang bab ini, $E, F$ adalah ruang bernorma atas $K = \R$ atau $\C$; dan
\emph{ruang Banach}\index{ruang Banach} berarti ruang bernorma
yang \omterm{def:b3:complete:complete}{lengkap}.

\section{Operator terbatas; ruang barisan}

\begin{definition}\label{def:b3:banach:operator}
Lambang $\mathcal L(E, F)$ menyatakan ruang berisi pemetaan linear yang
\emph{terbatas} (= \omterm{def:b3:topology:continuity}{kontinu}, Tahun ke-2) dengan \emph{norma
operator}\index{norma operator} $\vertiii T = \sup_{\norm x \leq
1}\norm{Tx}$; norma ini submultiplikatif: $\vertiii{ST} \leq
\vertiii S\,\vertiii T$. Adapun \emph{dual}\index{ruang dual} miliknya adalah
$E' = \mathcal L(E, K)$.
\end{definition}

\begin{proposition}\label{prop:b3:banach:llcomplete}
Jika $F$ ruang Banach, maka $\mathcal L(E, F)$ pun demikian;
khususnya $E'$ selalu merupakan ruang Banach.
\end{proposition}

\begin{proof}
Misalkan $(T_n)$ Cauchy terhadap $\vertiii\cdot$. Untuk setiap $x$ berlaku
$\norm{T_nx - T_mx} \leq \vertiii{T_n - T_m}\norm x$: sehingga $(T_nx)$
Cauchy di $F$, jadi konvergen; sebutlah limitnya $Tx$. Pemetaan $T$ bersifat
linear (sebagai limit identitas linear); lalu mengambil limit pada
$\norm{T_nx - T_mx} \leq \varepsilon\norm x$ (untuk $n, m \geq N$)
memberikan $\norm{T_nx - Tx} \leq \varepsilon\norm x$: jadi $T_n \to T$
dalam \omterm{def:b3:banach:operator}{norma operator}, dan $\vertiii T \leq \vertiii{T_N} +
\varepsilon < \infty$.
\end{proof}

\begin{definition}\label{def:b3:banach:lp}
Ruang barisan klasiknya\index{ruang barisan $\ell^p$}
(atas $K$, diindeks oleh $\N$):
\[
\ell^p = \Bigl\{x = (x_n) : \norm x_p = \Bigl(\sum_n
\abs{x_n}^p\Bigr)^{1/p} < \infty\Bigr\}\quad (1 \leq p <
\infty),
\]
\[
\ell^\infty = \{x : \norm x_\infty = \sup\abs{x_n} < \infty\},
\]
dan $c_0 = \{x : x_n \to 0\}$ dengan $\norm\cdot_\infty$. Bahwa
$\norm\cdot_p$ merupakan norma menyusul dari ketaksamaan Minkowski,
yang dibuktikan pada kasus diskretnya persis seperti di \cref{ch:b3:lp} (atau
dengan menjumlahkan ketaksamaan berdimensi berhingga pada Tahun ke-2). Semuanya
ruang Banach, dan $c_0$ merupakan subruang tertutup
$\ell^\infty$ (\cref{exo:b3:banach:3}).
\end{definition}

\begin{omfigure}
\begin{tikzpicture}[scale=1.5]
  \draw[->] (-1.45,0) -- (1.45,0) node[right,font=\small]
    {$x_1$};
  \draw[->] (0,-1.45) -- (0,1.45) node[above,font=\small]
    {$x_2$};
  % p = 1 diamond
  \draw[omDef, thick] (1,0) -- (0,1) -- (-1,0) -- (0,-1)
    -- cycle;
  % p = 2 circle
  \draw[omProp, thick] (0,0) circle (1);
  % p = 4 superellipse |x|^4 + |y|^4 = 1
  \draw[omThm, thick, smooth]
    plot[domain=0:90, samples=46]
    ({cos(\x)/(cos(\x)^4 + sin(\x)^4)^0.25},
     {sin(\x)/(cos(\x)^4 + sin(\x)^4)^0.25});
  \draw[omThm, thick, smooth]
    plot[domain=90:180, samples=46]
    ({cos(\x)/(cos(\x)^4 + sin(\x)^4)^0.25},
     {sin(\x)/(cos(\x)^4 + sin(\x)^4)^0.25});
  \draw[omThm, thick, smooth]
    plot[domain=180:270, samples=46]
    ({cos(\x)/(cos(\x)^4 + sin(\x)^4)^0.25},
     {sin(\x)/(cos(\x)^4 + sin(\x)^4)^0.25});
  \draw[omThm, thick, smooth]
    plot[domain=270:360, samples=46]
    ({cos(\x)/(cos(\x)^4 + sin(\x)^4)^0.25},
     {sin(\x)/(cos(\x)^4 + sin(\x)^4)^0.25});
  % p = infty square
  \draw[black!60, thick] (-1,-1) rectangle (1,1);
  \node[omDef, font=\small] at (0.40,0.40) {$p{=}1$};
  \node[omProp, font=\small] at (-0.52,0.75) {$p{=}2$};
  \node[omThm, font=\small] at (0.90,-0.52) {$p{=}4$};
  \node[black!60, font=\small] at (-1.62,1.15)
    {$p{=}\infty$};
  \draw[black!60, ->] (-1.32,1.11) -- (-1.02,1.0);
\end{tikzpicture}

{\small Bola satuan norma-$p$ di bidang, yang bersarang ketika
$p$ tumbuh dari $1$ (belah ketupat) melalui $2$ (cakram) dan $4$
(superelips) sampai $\infty$ (persegi). Kekonveksan setiap bolanya
\emph{merupakan} ketaksamaan Minkowski; sedangkan sudut-sudut di $p = 1$
dan $p = \infty$ adalah tempat kekonveksan sejati, ketunggalan
hampiran terbaik, dan kasus kesamaan pada
\cref{exo:b3:lp:12} merosot sekaligus.}
\end{omfigure}

\begin{proposition}[Deret Neumann]\label{prop:b3:banach:neumann}
Misalkan $E$ ruang Banach dan $T \in \mathcal L(E) = \mathcal L(E,E)$
dengan $\vertiii T < 1$. Maka $I - T$ terbalikkan di $\mathcal
L(E)$, dengan $(I - T)^{-1} = \sum_{n\geq0}T^n$ (yang konvergen dalam
\omterm{def:b3:banach:operator}{norma operator}). Akibatnya himpunan operator terbalikkan bersifat
terbuka, dan pembalikannya \omterm{def:b3:topology:continuity}{kontinu} di sana.\index{deret Neumann}
\end{proposition}

\begin{proof}
Deretnya konvergen mutlak (sebab $\vertiii{T^n} \leq \vertiii
T^n$, yang geometri) di dalam $\mathcal L(E)$ yang Banach
(\cref{prop:b3:banach:llcomplete};
\cref{exo:b3:complete:1}(b)). Secara teleskopik, $(I - T)\sum_{n \leq
N}T^n = I - T^{N+1} \to I$, dan serupa pada sisi yang lain.
Untuk keterbukaannya: jika $S$ terbalikkan dan $\vertiii{H} <
1/\vertiii{S^{-1}}$, maka $S + H = S(I + S^{-1}H)$ dengan
$\vertiii{S^{-1}H} < 1$: jadi terbalikkan. \omterm{def:b3:topology:continuity}{Kekontinuan} pembalikannya:
ungkapan deretnya memberikan $\vertiii{(S+H)^{-1} - S^{-1}} =
O(\vertiii H)$ secara lokal.
\end{proof}

\begin{example}[Sebuah persamaan Volterra yang diselesaikan lewat Neumann]
\label{ex:b3:banach:volterraexample}
Pada $E = \mathcal C(\intcc01)$, tinjaulah persamaan integral
\[
u(x) = 1 + \lambda\int_0^xu(t)\,\dd t,
\qquad\text{yakni}\qquad u = \mathbf 1 + \lambda Tu,
\quad (Tu)(x) = \int_0^xu .
\]
Di sini $\vertiii T \leq 1$, sehingga untuk $\abs\lambda < 1$ \omterm{prop:b3:banach:neumann}{deret Neumann}
langsung berlaku: $u = (I - \lambda T)^{-1}\mathbf 1 =
\sum_n\lambda^nT^n\mathbf 1$. Dengan menghitungnya, $T^n\mathbf 1 =
\frac{x^n}{n!}$, sehingga
\[
u(x) = \sum_{n\geq0}\frac{(\lambda x)^n}{n!} =
\eu^{\lambda x} ,
\]
seperti dibenarkan oleh penurunannya. Lebih baik lagi: $\vertiii{T^n} \leq
\frac1{n!}$ (sebab kernel beriterasinya menyusut secara faktorial), sehingga
$\sum\lambda^nT^n$ konvergen untuk \emph{setiap} $\lambda$ ---
jadi operator $I - \lambda T$ terbalikkan untuk setiap $\lambda
\in \C$, walaupun akhirnya $\vertiii{\lambda T} \geq 1$:
sebab yang menentukan adalah peluruhan spektral pangkatnya, bukan
norma pertamanya. Operator Volterra membawa peluruhan faktorial itu
sejak lahir (\cref{exo:b3:complete:4}(a) memanfaatkan persis
hal ini), dan itulah sebabnya masalah nilai awal tak pernah mengalami
gejala resonansi seperti masalah nilai batas
(\cref{ch:b3:spectral}).
\end{example}

\section{Hahn--Banach}

\begin{theorem}[Hahn--Banach, bentuk analitik]
\label{thm:b3:banach:hahnbanach}
Misalkan $E$ sebuah ruang vektor \emph{real}, $p \colon E \to \R$
bersifat \emph{sublinear} ($p(x + y) \leq p(x) + p(y)$ dan $p(tx) =
tp(x)$ untuk $t \geq 0$), $F \subseteq E$ sebuah subruang, dan $f
\colon F \to \R$ linear dengan $f \leq p$ pada $F$. Maka $f$ dapat diperluas
menjadi $\tilde f \colon E \to \R$ yang linear dengan $\tilde f \leq p$
pada $E$.\index{teorema Hahn--Banach}
\end{theorem}

\begin{proof}
\emph{Perluasan satu langkah.} Misalkan $x_0 \notin F$; kita perluas $f$ ke
$F \oplus \R x_0$ dengan memilih $\alpha = \tilde f(x_0)$
secara tepat: kita memerlukan, untuk setiap $y \in F$ dan $t > 0$,
\[
f(y) + t\alpha \leq p(y + tx_0)
\quad\text{dan}\quad
f(y) - t\alpha \leq p(y - tx_0),
\]
yang setelah dibagi $t$ (lewat kesublinearannya) menyusut menjadi
\[
\sup_{v \in F}\ \bigl[f(v) - p(v - x_0)\bigr]
\;\leq\; \alpha \;\leq\;
\inf_{u \in F}\ \bigl[p(u + x_0) - f(u)\bigr].
\]
Bilangan $\alpha$ semacam itu ada jika dan hanya jika setiap anggota kirinya $\leq$ setiap
anggota kanannya: dan memang $f(v) + f(u) = f(u + v) \leq p(u + v) \leq
p(u + x_0) + p(v - x_0)$, yakni $f(v) - p(v - x_0) \leq p(u +
x_0) - f(u)$.

\emph{Zorn.} Urutkan semua perluasan $f$ yang didominasi $p$
(berupa pasangan: subruang dan fungsional) menurut perluasannya; sebuah rantai
mempunyai gabungannya sebagai batas atas; dan unsur maksimalnya haruslah terdefinisi pada seluruh
$E$, sebab kalau tidak perluasan satu langkahnya akan bertentangan dengan kemaksimalannya.
\end{proof}

\begin{corollary}\label{cor:b3:banach:hbnormed}
Misalkan $E$ sebuah ruang bernorma (dengan $K = \R$ atau $\C$).
\begin{enumerate}
  \item Setiap $f \in F'$ (dengan $F$ sebuah subruang) dapat diperluas menjadi $\tilde f
        \in E'$ dengan $\norm{\tilde f}_{E'} = \norm f_{F'}$.
  \item Untuk setiap $x \neq 0$ ada $f \in E'$ dengan $\norm f
        = 1$ dan $f(x) = \norm x$. Khususnya $E'$ memisahkan
        titik $E$, dan $\norm x = \sup_{\norm f \leq
        1}\abs{f(x)}$.
  \item Untuk subruang tertutup $F$ dan $x \notin F$, ada $f
        \in E'$ yang lenyap pada $F$ dengan $f(x) = d(x, F)$ dan
        $\norm f \leq 1$.
\end{enumerate}
\end{corollary}

\begin{proof}
(1) Kasus realnya: terapkan \cref{thm:b3:banach:hahnbanach} dengan $p(x)
= \norm f_{F'}\,\norm x$ (yang sublinear); maka perluasannya memenuhi
$\pm\tilde f(x) = \tilde f(\pm x) \leq p(x)$, sehingga
$\norm{\tilde f} \leq \norm f$, sedangkan $\geq$ berasal dari pembatasannya.
Kasus kompleksnya: misalkan $u = \operatorname{Re}f$, yaitu fungsional real
dengan $\abs u \leq \norm f\norm\cdot$; catat bahwa $f(x) = u(x) -
\iu\,u(\iu x)$ (periksa pada bagian real dan imajinernya:
$\operatorname{Im}f(x) = -\operatorname{Re}f(\iu x)$). Perluaslah
$u$ secara linear-real dengan batas yang sama, lalu tetapkan $\tilde f(x) =
\tilde u(x) - \iu\tilde u(\iu x)$: ia linear atas $\C$ (periksa langsung pada
perkalian dengan $\iu$) dan memperluas $f$; normanya: untuk $x$ yang diberikan
tulis $\tilde f(x) = r\eu^{\iu\theta}$, maka $\abs{\tilde f(x)}
= \tilde f(\eu^{-\iu\theta}x) = \tilde u(\eu^{-\iu\theta}x)
\leq \norm f\,\norm x$.

(2) Pada $F = Kx$ definisikan $f(tx) = t\norm x$: yang bernorma $1$ pada $F$;
lalu perluas lewat (1). Untuk rumus dualitasnya: $\leq$ jelas, sedangkan $\geq$ lewat
$f$ ini.

(3) Pada $F \oplus Kx$ definisikan $f(y + tx) = t\,d(x, F)$; maka untuk
$t \neq 0$ berlaku $\norm{y + tx} = \abs t\,\norm{x + y/t} \geq \abs
t\,d(x, F) = \abs{f(y + tx)}$: sehingga $\norm f \leq 1$ pada
subruangnya; lalu perluas lewat (1).
\end{proof}

\begin{remark}\label{rem:b3:banach:bidual}
Menurut (2), pemetaan kanonik $J \colon E \to E''$ dengan $J(x)(f) =
f(x)$ merupakan isometri (\cref{exo:b3:banach:10}): jadi setiap ruang
bernorma duduk di dalam bidualnya. Ruang yang $J$-nya surjektif
disebut \emph{refleksif}\index{ruang refleksif}; dan soal akhir
pekan menunjukkan bahwa $\ell^p$ (untuk $1 < p < \infty$) \omterm{rem:b3:banach:bidual}{refleksif} sedangkan
$\ell^1$ tidak.
\end{remark}

\section{Trilogi Baire}

\begin{theorem}[Banach--Steinhaus, keterbatasan seragam]
\label{thm:b3:banach:banachsteinhaus}
Misalkan $E$ sebuah ruang \emph{Banach}, $F$ bernorma, dan $(T_i)_{i\in
I} \subseteq \mathcal L(E, F)$ sebuah keluarga dengan $\sup_i
\norm{T_ix} < \infty$ untuk setiap $x \in E$. Maka $\sup_i
\vertiii{T_i} < \infty$.\index{teorema Banach--Steinhaus}
\end{theorem}

\begin{proof}
Himpunan $F_n = \{x : \sup_i\norm{T_ix} \leq n\}$ bersifat tertutup
(sebagai irisan prapeta bola tertutup) dan menyelimuti $E$.
Lalu Baire (\cref{thm:b3:complete:baire}) memberikan $n_0$ dan sebuah bola
$B(x_0, r) \subseteq F_{n_0}$. Untuk $\norm z < r$: $\norm{T_iz}
\leq \norm{T_i(x_0 + z)} + \norm{T_ix_0} \leq 2n_0$, sehingga
$\vertiii{T_i} \leq 2n_0/r$ untuk setiap $i$.
\end{proof}

\begin{corollary}\label{cor:b3:banach:bslimits}
Jika $E$ Banach dan $T_n \in \mathcal L(E,F)$ konvergen
\emph{titik demi titik} (yakni $T_nx \to Tx$ untuk setiap $x$), maka $\sup_n
\vertiii{T_n} < \infty$, $T \in \mathcal L(E, F)$, dan
$\vertiii T \leq \liminf \vertiii{T_n}$.
\end{corollary}

\begin{proof}
Barisan konvergen bersifat terbatas: itulah keterbatasan titik demi titiknya;
lalu Banach--Steinhaus membatasi normanya oleh suatu $M$; sehingga $\norm{Tx}
= \lim\norm{T_nx} \leq M\norm x$ (sebab $T$ linear sebagai limit titik demi
titik), dan batas yang lebih tajamnya lewat mengambil $\liminf$ pada
$\norm{T_nx} \leq \vertiii{T_n}\norm x$.
\end{proof}

\begin{theorem}[Deret Fourier yang divergen]
\label{thm:b3:banach:fourierdiverge}
Ada fungsi \omterm{def:b3:topology:continuity}{kontinu} berperiode $2\pi$, sebut $f$, yang
deret Fouriernya divergen di $0$: $\sup_N\abs{S_N(f)(0)} =
\infty$. Bahkan $f$ semacam itu membentuk himpunan bagian padat di $\mathcal
C(S^1)$.\index{deret Fourier, divergensi}
\end{theorem}

\begin{proof}
Bekerjalah di $E = (\mathcal C(S^1), \norm\cdot_\infty)$, sebuah ruang
Banach, dengan fungsional $\Lambda_N(f) = S_N(f)(0) =
\frac1{2\pi}\int_{-\pi}^{\pi}f(t)\,D_N(t)\,\dd t$ (yaitu kernel
Dirichlet, Tahun ke-2). Setiap $\Lambda_N$ \omterm{def:b3:topology:continuity}{kontinu} dengan
\[
\norm{\Lambda_N} = \frac1{2\pi}\int_{-\pi}^\pi\abs{D_N(t)}\,\dd
t \;=\; L_N .
\]
($\leq$ jelas; sedangkan $\geq$: ambillah $f$ yang \omterm{def:b3:topology:continuity}{kontinu} dengan $\norm f_\infty
\leq 1$ yang menghampiri $\operatorname{sign}D_N$ --- sebab tandanya
berlompatan berhingga banyak; dan menghaluskan setiap lompatannya pada selang
berpanjang $\varepsilon$ mengubah integralnya sebesar $O(N\varepsilon)$.)
\emph{Konstanta Lebesgue} $L_N$ menuju tak hingga:
\[
L_N = \frac1{2\pi}\int_{-\pi}^{\pi}
\frac{\abs{\sin\bigl((N{+}\tfrac12)t\bigr)}}{\abs{\sin(t/2)}}
\,\dd t
\geq \frac{2}{\pi}\int_0^{\pi}
\frac{\abs{\sin\bigl((N{+}\tfrac12)t\bigr)}}{t}\,\dd t
= \frac{2}{\pi}\int_0^{(N+\frac12)\pi}\frac{\abs{\sin
u}}{u}\,\dd u
\]
(dengan memakai $\abs{\sin(t/2)} \leq t/2$ pada $[0, \pi]$, lalu
menyubstitusikan $u = (N + \tfrac12)t$). Dengan memotongnya atas lengkungan:
\[
\int_{(k-1)\pi}^{k\pi}\frac{\abs{\sin u}}u\,\dd u
\geq \frac1{k\pi}\int_{(k-1)\pi}^{k\pi}\abs{\sin u}\,\dd u =
\frac{2}{k\pi},
\qquad\text{sehingga}\qquad
L_N \geq \frac{4}{\pi^2}\sum_{k=1}^N\frac1k
\xrightarrow[N\to\infty]{} \infty .
\]
Seandainya setiap $f$ yang \omterm{def:b3:topology:continuity}{kontinu} mempunyai $\sup_N
\abs{\Lambda_N(f)} < \infty$, maka Banach--Steinhaus akan memaksa
$\sup_N\norm{\Lambda_N} < \infty$: kontradiksi. Jadi ada $f$
--- bahkan himpunan $f$ yang padat dan tak \omterm{rem:b3:complete:meagre}{kurus} (yakni komplemen
$\bigcup_M\{f: \sup_N\abs{\Lambda_Nf}\leq M\}$, yaitu gabungan terbilang
himpunan tertutup yang, karena tak berinterior menurut uraian di atas yang
diterapkan pada bola mana pun, bersifat \omterm{rem:b3:complete:meagre}{kurus}) --- yang mempunyai $\sup_N\abs{S_Nf(0)} =
\infty$.
\end{proof}

\begin{theorem}[Pemetaan terbuka]\label{thm:b3:banach:openmapping}
Misalkan $E, F$ ruang Banach dan $T \in \mathcal L(E, F)$
bersifat \emph{surjektif}. Maka $T$ terbuka: $T(B_E(0,1)) \supseteq
B_F(0, c)$ untuk suatu $c > 0$. Akibatnya operator terbatas yang
bijektif antara ruang Banach berinvers
terbatas.\index{teorema pemetaan terbuka}
\end{theorem}

\begin{proof}
Tulis $B = B_E(0,1)$. Kesurjektifannya memberikan $F = \bigcup_n
\overline{T(nB)} = \bigcup_n n\,\overline{T(B)}$; lalu Baire
(\cref{thm:b3:complete:baire}) memberikan \omterm{def:b3:topology:interior}{interior} kepada
$\overline{T(B)}$: yakni ada $B_F(y_0, 4c) \subseteq
\overline{T(B)}$. Pusatkan ulang di $0$: untuk $\norm y < 4c$, baik
$y_0 + y$ maupun $y_0$ merupakan limit peta $Tu_k$, $Tv_k$ dengan
$u_k, v_k \in B$, sehingga $y = \lim T(u_k - v_k)$ dengan $u_k - v_k
\in 2B$: jadi $B_F(0, 4c) \subseteq \overline{T(2B)}$, yakni
$B_F(0, 2c) \subseteq \overline{T(B)}$.

\emph{Membuang \omterm{def:b3:topology:interior}{penutupnya}} (dan di sinilah \omterm{def:b3:complete:complete}{kelengkapan} $E$ masuk):
misalkan $\norm y < c$. Ambillah $x_1 \in \frac12 B$ dengan $\norm{y -
Tx_1} < c/2$ (sebab $B_F(0,2c) \subseteq \overline{T(B)}$ diskalakan
$\frac12$); lalu secara induktif $x_k \in 2^{-k}B$ dengan $\norm{y -
T(x_1 + \dots + x_k)} < c\,2^{-k}$. Deret $\sum x_k$
konvergen mutlak di $E$ yang Banach, ke $x \in B$ (bernorma $<
\sum 2^{-k} = 1$), dan $Tx = y$ menurut \omterm{def:b3:topology:continuity}{kekontinuannya}: jadi $B_F(0, c)
\subseteq T(B)$. Keterbukaan $T$ pada \omterm{def:b3:topology:topology}{himpunan terbuka} sembarang menyusul lewat
translasi dan penskalaan; sedangkan untuk akibatnya, keterbukaan $T$
berarti $T^{-1}$ \omterm{def:b3:topology:continuity}{kontinu}.
\end{proof}

\begin{corollary}[Norma yang ekuivalen]\label{cor:b3:banach:equivnorms}
Jika sebuah ruang vektor \omterm{def:b3:complete:complete}{lengkap} terhadap dua norma yang sebanding
($\norm\cdot_a \leq C\norm\cdot_b$), maka kedua normanya ekuivalen.
\end{corollary}

\begin{proof}
Identitas $(E, \norm\cdot_b) \to (E, \norm\cdot_a)$ bersifat
terbatas dan bijektif antara ruang Banach: sehingga inversnya
terbatas.
\end{proof}

\begin{theorem}[Grafik tertutup]\label{thm:b3:banach:closedgraph}
Misalkan $E, F$ Banach dan $T \colon E \to F$ linear. Jika
grafiknya $\Gamma = \{(x, Tx)\}$ tertutup di $E \times F$ (yakni
$x_n \to x$ dan $Tx_n \to y$ mengakibatkan $y = Tx$), maka $T$
terbatas.\index{teorema grafik tertutup}
\end{theorem}

\begin{proof}
Ruang $E \times F$ dengan $\norm{(x,y)} = \norm x + \norm y$ bersifat Banach;
dan $\Gamma$, sebagai subruang tertutup, juga Banach. Proyeksi
$\pi_E\colon \Gamma \to E$ bersifat terbatas dan bijektif, sehingga
inversnya $x \mapsto (x, Tx)$ terbatas
(\cref{thm:b3:banach:openmapping}): jadi $\norm{Tx} \leq \norm{(x,
Tx)} \leq C\norm x$.
\end{proof}

\begin{method}\label{met:b3:banach:trilogy}
Kapan harus meraih teorema yang mana. \emph{Hahn--Banach}: untuk
menghasilkan fungsional dengan perilaku yang ditentukan (menormakan sebuah
vektor, melenyap pada subruang, memperluas dari subruang) ---
dan tak perlu \omterm{def:b3:complete:complete}{kelengkapan}. \emph{Banach--Steinhaus}: untuk mengubah
informasi titik demi titik menjadi batas seragam --- biasanya untuk menunjukkan
bahwa operasi limitnya \omterm{def:b3:topology:continuity}{kontinu}, atau (lewat kontraposisi) untuk membuktikan
\emph{divergensi} bagi suatu unsur, seperti pada deret Fourier.
\emph{Pemetaan terbuka / grafik tertutup}: untuk memperoleh \omterm{def:b3:topology:continuity}{kekontinuan} secara cuma-cuma
dari kebijektifan aljabar atau dari sifat ketertutupan
grafiknya --- pemakaian khasnya: membandingkan dua norma \omterm{def:b3:complete:complete}{lengkap}, atau membuktikan
\omterm{def:b3:topology:continuity}{kekontinuan} otomatis. Ketiga teorema Baire itu menuntut
\omterm{def:b3:complete:complete}{kelengkapan} \emph{pada sumbernya}; dan tanpa itu ada contoh penyangkalnya
(\cref{exo:b3:banach:7}).
\end{method}

\section{Ruang dual, secara konkret}

\begin{theorem}\label{thm:b3:banach:duals}
Secara isometrik: $(c_0)' \cong \ell^1$ dan $(\ell^1)' \cong
\ell^\infty$, lewat pemasangan $\langle x, y\rangle = \sum_n
x_ny_n$.\index{dual ruang barisan}
\end{theorem}

\begin{proof}
Kita buktikan $(c_0)' \cong \ell^1$; sedangkan pengenalan keduanya adalah
\cref{exo:b3:banach:5}. Kepada $y \in \ell^1$ kaitkan $\Lambda_y(x)
= \sum x_ny_n$ (dengan $x \in c_0$): yang konvergen mutlak, dengan
$\abs{\Lambda_y(x)} \leq \norm x_\infty\norm y_1$, sehingga
$\norm{\Lambda_y} \leq \norm y_1$. Sebaliknya misalkan $\Lambda \in
(c_0)'$; tetapkan $y_n = \Lambda(e_n)$ (dengan $e_n$ barisan satuannya).
Untuk sembarang $N$, ujilah $x^{(N)} = \sum_{n \leq N}
\operatorname{sign}(\overline{y_n})\,e_n \in c_0$ (yang bernorma $\leq
1$; pada kasus kompleksnya pakailah faktor unimodular $\bar y_n/\abs
{y_n}$): maka $\Lambda(x^{(N)}) = \sum_{n\leq N}\abs{y_n} \leq
\norm\Lambda$. Jadi $y \in \ell^1$ dengan $\norm y_1 \leq
\norm\Lambda$. Akhirnya $\Lambda = \Lambda_y$: sebab keduanya bersesuaian pada
$e_n$, sehingga pada barisan berhingga, yang padat di $c_0$ (lewat pemenggalan:
$\norm{x - \sum_{n \leq N}x_ne_n}_\infty = \sup_{n > N}\abs{x_n}
\to 0$ persis karena $x_n \to 0$); dan fungsional \omterm{def:b3:topology:continuity}{kontinu}
yang bersesuaian pada \omterm{def:b3:topology:interior}{himpunan padat} pastilah sama. Korespondensinya linear,
bijektif, dan isometrik (sebab $\norm{\Lambda_y} = \norm y_1$ dari
kedua ketaksamaannya).
\end{proof}

\section{Latihan}

\begin{exercise}[$\star$]\label{exo:b3:banach:1}
Hitunglah \omterm{def:b3:banach:operator}{norma operatornya}:
(a) pergeseran $S(x_1, x_2, \dots) = (0, x_1, x_2, \dots)$ dan
$S^*(x_1, x_2, \dots) = (x_2, x_3, \dots)$ pada $\ell^2$;
(b) operator perkalian $M_a x = (a_nx_n)$ pada $\ell^2$,
untuk $a \in \ell^\infty$;
(c) fungsional $\Lambda(f) = \int_0^{1/2}f -
\int_{1/2}^1f$ pada $\mathcal C(\intcc01)$ --- tunjukkan
$\norm\Lambda = 1$ dan bahwa normanya \emph{tidak tercapai}.
\end{exercise}

\begin{exercise}[$\star$]\label{exo:b3:banach:2}
Misalkan $E$ Banach, $T \in \mathcal L(E)$ terbalikkan, dan $S$
memenuhi $\vertiii{S - T} < 1/\vertiii{T^{-1}}$. Tunjukkan bahwa $S$
terbalikkan lalu taksirlah $\vertiii{S^{-1} - T^{-1}}$. Terapannya:
jika sistem linear $Tx = b$ \omterm{def:b3:groups:derived}{terselesaikan} dengan $T$ terbalikkan, maka
usikan $T$ yang cukup kecil tetap membuatnya \omterm{def:b3:groups:derived}{terselesaikan} secara
tunggal, dengan batas kuantitatif pada perubahan penyelesaiannya.
\end{exercise}

\begin{exercise}[$\star\star$]\label{exo:b3:banach:3}
(a) Buktikan bahwa $\ell^1$, $\ell^\infty$ dan $c_0$ merupakan ruang
Banach, dan bahwa $c_0$ adalah \omterm{def:b3:topology:interior}{penutup} di $\ell^\infty$ dari
ruang berisi barisan berhingga.
(b) Tunjukkan $\ell^p \subseteq \ell^q$ dengan $\norm\cdot_q \leq
\norm\cdot_p$ untuk $1 \leq p \leq q \leq \infty$, dan bahwa
inklusinya sejati.
\end{exercise}

\begin{exercise}[$\star\star$]\label{exo:b3:banach:4}
Misalkan $F \subseteq E$ sebuah subruang tertutup dan $x \notin F$.
Dengan memakai \cref{cor:b3:banach:hbnormed}, buktikan rumus dualitas
\[
d(x, F) = \max\bigl\{\abs{f(x)} : f \in E',\ \norm f \leq 1,\
f\restriction_F = 0\bigr\}
\]
(perhatikan: sebuah \emph{maksimum}). Simpulkan bahwa $F = \bigcap\{\ker f : f
\in E',\ f\restriction_F = 0\}$: jadi subruang tertutup tepat berupa
irisan kernel fungsional.
\end{exercise}

\begin{exercise}[$\star\star$]\label{exo:b3:banach:5}
Buktikan $(\ell^1)' \cong \ell^\infty$ secara isometrik, mengikuti
skema \cref{thm:b3:banach:duals} (barisan berhingga padat
di $\ell^1$).
Di manakah argumennya runtuh untuk $(\ell^\infty)'$?
\end{exercise}

\begin{exercise}[$\star\star$]\label{exo:b3:banach:6}
Misalkan $E, F, G$ bernorma dengan $E$ Banach, dan $B \colon E \times
F \to G$ bilinear serta \omterm{def:b3:topology:continuity}{kontinu} pada setiap variabelnya secara terpisah. Tunjukkan
bahwa $B$ \omterm{def:b3:topology:continuity}{kontinu} (secara bersama): $\norm{B(x,y)} \leq
C\norm x\norm y$. \emph{(Terapkan Banach--Steinhaus pada keluarga
$(B(\cdot, y))_{\norm y \leq 1}$.)}
\end{exercise}

\begin{exercise}[$\star\star$]\label{exo:b3:banach:7}
(a) Pada $E = \mathcal C(\intcc01)$, bandingkan $\norm\cdot_\infty$
dan $\norm\cdot_1$: identitas $(E, \norm\cdot_\infty) \to
(E, \norm\cdot_1)$ bersifat terbatas dan bijektif tetapi inversnya
tak terbatas. Hipotesis mana pada
\cref{cor:b3:banach:equivnorms} yang gagal?
(b) Tampilkan pemetaan linear tak \omterm{def:b3:topology:continuity}{kontinu} dari sebuah subruang padat
$\ell^2$ ke $K$ (misalnya pada barisan berhingga), lalu jelaskan mengapa
hal itu tak bertentangan dengan teorema grafik tertutup.
\end{exercise}

\begin{exercise}[$\star\star\star$]\label{exo:b3:banach:8}
(Hellinger--Toeplitz) Misalkan $T \colon \ell^2 \to \ell^2$ linear
(terdefinisi di mana-mana) dan \emph{simetris}: $\langle Tx,
y\rangle = \langle x, Ty\rangle$ untuk setiap $x, y$, dengan $\langle
x, y \rangle = \sum x_n\bar y_n$. Tunjukkan bahwa $T$ terbatas.
\emph{(Lewat grafik tertutup: jika $x_k \to x$ dan $Tx_k \to z$, ujilah
terhadap $y$ sembarang.)} Pelajarannya: operator simetris yang tak terbatas
--- yakni Hamiltonian pada mekanika kuantum --- tak pernah dapat
didefinisikan pada seluruh ruangnya.
\end{exercise}

\begin{exercise}[$\star\star\star$]\label{exo:b3:banach:9}
(Teorema Polya tentang kuadratur) Untuk setiap $n$, misalkan $\Lambda_n(f)
= \sum_{i=0}^{n} w_{i,n}f(x_{i,n})$ sebuah aturan kuadratur pada
$\mathcal C(\intcc01)$ (dengan $x_{i,n} \in \intcc01$, $w_{i,n} \in
\R$). Tunjukkan bahwa $\Lambda_n(f) \to \int_0^1f$ untuk \emph{setiap}
$f$ yang \omterm{def:b3:topology:continuity}{kontinu} jika dan hanya jika: (i) $\Lambda_n(P) \to \int_0^1P$
untuk setiap polinomial $P$, dan (ii) $\sup_n\sum_i\abs{w_{i,n}} <
\infty$. \emph{(Hitunglah $\norm{\Lambda_n}$; lalu pakai
Banach--Steinhaus dan Weierstrass.)} Periksalah bahwa aturan dengan
bobot \emph{positif} yang eksak pada konstanta memenuhi (ii)
secara otomatis.
\end{exercise}

\begin{exercise}[$\star\star$]\label{exo:b3:banach:10}
Tunjukkan bahwa $J \colon E \to E''$ dengan $J(x)(f) = f(x)$ merupakan
isometri linear (pakailah \cref{cor:b3:banach:hbnormed}(2)), dan bahwa ia
surjektif bila $\dim E < \infty$. Tunjukkan pula bahwa jika $E'$
\omterm{def:b3:galois:separable}{separabel} maka $E$ pun \omterm{def:b3:galois:separable}{separabel}. \emph{(Ambil $x_n$ yang hampir menormakan
barisan padat di $E'$ lalu tunjukkan bahwa rentang tertutupnya adalah $E$, lewat
\cref{cor:b3:banach:hbnormed}(3).)}
\end{exercise}

\begin{exercise}[$\star\star$]\label{exo:b3:banach:11}
(Ruang kuosien) Misalkan $E$ sebuah ruang Banach dan $F \subseteq
E$ sebuah subruang \emph{tertutup}. Pada $E/F$ definisikan
\[
\norm{\bar x} \;=\; d(x, F) = \inf_{y\in F}\norm{x - y} .
\]
(a) Tunjukkan bahwa ini merupakan \emph{norma} yang terdefinisi dengan baik pada $E/F$ (di
manakah ketertutupan $F$ masuk?), dan bahwa proyeksi $\pi
\colon E \to E/F$ mempunyai $\vertiii\pi \leq 1$ serta memetakan bola
satuan terbuka pada bola satuan terbuka.
(b) Tunjukkan bahwa $E/F$ \omterm{def:b3:complete:complete}{lengkap}. \emph{(Pakailah kriteria deret
pada \cref{exo:b3:complete:1}(b): diberikan kelas
$\bar x_k$ dengan $\sum\norm{\bar x_k} < \infty$, angkatlah masing-masing ke
$x_k \in E$ dengan $\norm{x_k} \leq \norm{\bar x_k} + 2^{-k}$
lalu jumlahkan di $E$.)}
(c) Hitunglah: untuk $E = c$ (barisan konvergen) dan $F =
c_0$, tunjukkan $c/c_0 \cong K$ secara isometrik lewat $\bar x \mapsto
\lim_nx_n$.
\end{exercise}

\begin{exercise}[$\star\star$]\label{exo:b3:banach:12}
(Proyeksi terbatas dan subruang berkomplemen) Misalkan $E$ sebuah
ruang Banach dan $P \colon E \to E$ linear dengan $P^2 = P$
(yakni proyeksi aljabar), $V = \operatorname{im}P$, dan $W =
\ker P$.
(a) Andaikan $P$ terbatas. Tunjukkan bahwa $V$ dan $W$ tertutup
dan $E = V \oplus W$ dengan penguraian $x = Px + (x -
Px)$.
(b) Sebaliknya, andaikan $E = V \oplus W$ dengan \emph{keduanya},
yakni $V, W$, tertutup, lalu misalkan $P$ proyeksi pada $V$ sepanjang $W$.
Tunjukkan bahwa $P$ terbatas. \emph{(Lewat grafik tertutup: jika $x_n \to x$
dan $Px_n \to z$, maka $z \in V$, $x_n - Px_n \to x - z \in
W$, dan ketunggalan penguraiannya mengenali $z =
Px$.)}
(c) Simpulkan \emph{kesetaraannya}: sebuah subruang $V$ mempunyai
proyeksi terbatas jika dan hanya jika ia tertutup dan mempunyai komplemen
aljabar yang tertutup --- lalu catat (tanpa bukti) bahwa
ada subruang tertutup tanpa sifat itu ($c_0$ di dalam
$\ell^\infty$ merupakan contoh klasiknya): jadi ruang Hilbert,
tempat $V^\perp$ selalu berhasil
(\cref{ch:b3:hilbert}), merupakan perkecualian, bukan aturannya.
\end{exercise}

\section{Soal: dualitas ruang \texorpdfstring{$\ell^p$}%
{lp}}

\begin{problem}[{Soal akhir pekan --- $(\ell^p)' = \ell^q$,
kerefleksifan, dan keanehan $\ell^\infty$}]
\label{pb:b3:banach:1}
Tetapkan $1 < p < \infty$ dan misalkan $q$ eksponen sekawannya, yakni
$\frac1p + \frac1q = 1$. Pemasangannya sepanjang soal ini adalah $\langle x,
y\rangle = \sum_n x_ny_n$.

\medskip
\noindent\textbf{Bagian I --- Hölder dan Minkowski untuk
barisan.}
\begin{enumerate}
  \item (Ketaksamaan Young) Untuk $a, b \geq 0$ tunjukkan $ab \leq
        \frac{a^p}p + \frac{b^q}q$, dengan memakai kekonkafan $\log$ atau
        dengan mempelajari $t \mapsto \frac{t^p}p + \frac1q - t$.
  \item (Hölder) Simpulkan: $\abs{\langle x, y\rangle} \leq
        \norm x_p\norm y_q$ untuk $x \in \ell^p$, $y \in \ell^q$;
        lalu kenali kasus kesamaannya.
  \item (Minkowski) Simpulkan ketaksamaan segitiga bagi
        $\norm\cdot_p$. \emph{(Tulis $\abs{x_n + y_n}^p \leq
        \abs{x_n}\,\abs{x_n{+}y_n}^{p-1} +
        \abs{y_n}\,\abs{x_n{+}y_n}^{p-1}$ lalu terapkan Hölder pada
        setiap sukunya.)}
  \item Buktikan bahwa $\ell^p$ \omterm{def:b3:complete:complete}{lengkap} dan bahwa barisan
        berhingga padat di dalamnya.
\end{enumerate}

\medskip
\noindent\textbf{Bagian II --- Dualitas $(\ell^p)' =
\ell^q$.}
\begin{enumerate}[resume]
  \item Untuk $y \in \ell^q$, tunjukkan bahwa $\Lambda_y(x) = \langle
        x, y\rangle$ mendefinisikan $\Lambda_y \in (\ell^p)'$ dengan
        $\norm{\Lambda_y} \leq \norm y_q$, lalu, dengan mengujinya pada
        $x_n = \abs{y_n}^{q-1}\operatorname{sign}(y_n)$
        (yang dipenggal dan dinormalkan secukupnya), bahwa
        $\norm{\Lambda_y} = \norm y_q$.
  \item Sebaliknya, diberikan $\Lambda \in (\ell^p)'$, tetapkan $y_n =
        \Lambda(e_n)$; tunjukkan $y \in \ell^q$ dengan $\norm y_q \leq
        \norm\Lambda$ \emph{(ujilah pada pemenggalannya seperti pada pertanyaan
        5 lalu biarkan panjang pemenggalannya tumbuh)}, dan simpulkan
        $\Lambda = \Lambda_y$: jadi pemetaan $y \mapsto \Lambda_y$ merupakan
        isomorfisma isometrik $\ell^q \to (\ell^p)'$.
  \item Simpulkan bahwa $\ell^p$ bersifat \emph{\omterm{rem:b3:banach:bidual}{refleksif}} untuk $1 < p <
        \infty$: sebab dengan mengomposisikan kedua dualitasnya, setiap unsur
        $(\ell^p)''$ berasal dari $\ell^p$; lalu periksalah dengan cermat bahwa
        komposisinya memang $J$ yang kanonik.
\end{enumerate}

\medskip
\noindent\textbf{Bagian III --- $\ell^1$ dan $\ell^\infty$ adalah
hewan yang berbeda.}
\begin{enumerate}[resume]
  \item Tunjukkan bahwa $\ell^p$ (untuk $1 \leq p < \infty$) dan $c_0$ bersifat
        \omterm{def:b3:galois:separable}{separabel}, tetapi $\ell^\infty$ tidak. \emph{(Sebab
        barisan indikator dari himpunan bagian $\N$ yang tak terbilang banyaknya
        saling berjarak $1$ berpasangan.)}
  \item Simpulkan dari \cref{exo:b3:banach:10} bahwa $(\ell^1)'
        \cong \ell^\infty$ tetapi $(\ell^\infty)' \not\cong
        \ell^1$: jadi $\ell^1$ \emph{tidak} \omterm{rem:b3:banach:bidual}{refleksif}.
        \emph{(Sebab jika $(\ell^\infty)'$ adalah $\ell^1$, ia akan
        \omterm{def:b3:galois:separable}{separabel}, sehingga memaksa $\ell^\infty$ \omterm{def:b3:galois:separable}{separabel}.)}
  \item (Limit Banach, secara eksplisit) Pada $\ell^\infty_\R$, misalkan
        $p(x) = \limsup_n \frac{x_1 + \dots + x_n}{n}$. Tunjukkan
        bahwa $p$ sublinear, dan bahwa pada subruang $c$ berisi
        barisan konvergen, $\mathrm{LIM}(x) = \lim x$
        memenuhi $\mathrm{LIM} \leq p$. Perluaslah lewat Hahn--Banach
        menjadi $\mathrm{LIM} \colon \ell^\infty_\R \to \R$ lalu tunjukkan:
        $\mathrm{LIM}$ bersifat positif (yakni $x \geq 0 \Rightarrow
        \mathrm{LIM}(x) \geq 0$), invarian terhadap pergeseran
        (yakni $\mathrm{LIM}(x_2, x_3, \dots) = \mathrm{LIM}(x)$),
        memperluas limitnya, serta memenuhi $\liminf x \leq
        \mathrm{LIM}(x)\leq \limsup x$.
  \item Tunjukkan bahwa $\mathrm{LIM}$ semacam itu, dipandang di
        $(\ell^\infty)'$, \emph{bukan} berbentuk $\Lambda_y$
        untuk $y \in \ell^1$ mana pun; lalu simpulkan lagi $(\ell^\infty)'
        \neq \ell^1$. \emph{(Nilaikan pada barisan satuan
        $e_n$, lalu pada barisan konstan $1$.)}
  \item Nilaikan $\mathrm{LIM}$ pada $(0,1,0,1,\dots)$, lalu tunjukkan
        bahwa tak ada perluasan limit yang \emph{multiplikatif} dan
        invarian terhadap pergeseran \emph{(tinjaulah $x =
        (0,1,0,1,\dots)$ dan $x\cdot Sx$ dengan $S$ berupa
        pergeserannya)}.
\end{enumerate}

\medskip
\noindent\textbf{Bagian IV --- \omterm{def:b3:topology:interior}{Penutup}: mengapa kerefleksifan
penting.}
\begin{enumerate}[resume]
  \item Dengan memakai \cref{cor:b3:banach:bslimits} dan pertanyaan 6,
        tunjukkan bahwa setiap barisan terbatas di $\ell^p$ (untuk $1 < p <
        \infty$) mempunyai subbarisan $(x^{(k)})$ yang konvergen
        \emph{secara lemah}: yakni $\Lambda(x^{(k)})$ konvergen untuk setiap
        $\Lambda \in (\ell^p)'$. \emph{(Lewat ekstraksi diagonal pada
        koordinatnya yang terbilang; lalu kenali limit lemahnya
        di $\ell^p$ dengan memakai keterbatasan seragam normanya dan
        Hölder.)} Tunjukkan lewat contoh ($e_n$ di $\ell^1$, terhadap
        unsur $\ell^\infty$ yang dipilih baik-baik) bahwa hal ini gagal
        di $\ell^1$: jadi kekompakan lemah merupakan hak istimewa \omterm{rem:b3:banach:bidual}{ruang
        refleksif}.
\end{enumerate}

\medskip
\noindent\textbf{Bagian V --- \omterm{def:b3:topology:topology}{Topologi} lemah yang bekerja, dan
kejutan Schur.} Tulis $x^{(k)} \rightharpoonup x$ di dalam sebuah
ruang bernorma $E$ (yakni \emph{kekonvergenan lemah}) bila
$\Lambda(x^{(k)}) \to \Lambda(x)$ untuk setiap $\Lambda \in E'$.
\begin{enumerate}[resume]
  \item Lengkapilah sensusnya: tunjukkan $(c_0)' \cong \ell^1$
        secara isometrik, lewat skema pertanyaan 5--6 (apa yang
        menggantikan barisan pengujinya?). Lalu rangkailah rantai
        $c_0 \to \ell^1 \to \ell^\infty \to \dots$ berisi \omterm{def:b3:banach:operator}{dual}
        berturutan dan tandai di mana kerefleksifannya gagal.
  \item Tunjukkan bahwa setiap barisan yang konvergen lemah di ruang Banach
        bersifat terbatas: pandanglah $x^{(k)}$ lewat
        pembenaman kanonik $J$ sebagai fungsional pada $E'$ lalu
        terapkan Banach--Steinhaus
        (\cref{thm:b3:banach:banachsteinhaus}) --- pada ruang
        Banach yang mana, dan mengapa \omterm{def:b3:complete:complete}{kelengkapannya} tersedia di sana?
  \item Tunjukkan bahwa di $\ell^p$ dengan $1 < p < \infty$: $x^{(k)}
        \rightharpoonup x$ jika dan hanya jika $\sup_k\norm{x^{(k)}}_p <
        \infty$ dan $x^{(k)}_n \to x_n$ untuk setiap koordinat
        $n$ \emph{(satu arahnya memakai Banach--Steinhaus lewat
        pembenaman kanoniknya; sedangkan untuk arah lainnya, hampiri $y
        \in \ell^q$ dengan barisan berhingga)}. Lalu simpulkan $e_k
        \rightharpoonup 0$ di $\ell^2$ padahal $\norm{e_k}_2 =
        1$: jadi limit lemah dapat kehilangan massa.
  \item Tunjukkan bahwa normanya semikontinu bawah secara lemah:
        yakni $x^{(k)} \rightharpoonup x$ mengakibatkan $\norm x \leq
        \liminf_k\,\norm{x^{(k)}}$ \emph{(ambillah fungsional
        penormaan bagi $x$, \cref{cor:b3:banach:hbnormed})}.
  \item (Radon--Riesz di $\ell^2$) Tunjukkan bahwa di $\ell^2$, kekonvergenan
        lemah bersama kekonvergenan normanya mengakibatkan kekonvergenan
        norma \emph{(jabarkan $\norm{x^{(k)} -
        x}_2^2$)}. Berilah contoh penyangkal bagi pernyataan yang sama
        tanpa hipotesis normanya.
  \item (Schur, langkah 1) Misalkan $x^{(k)} \rightharpoonup 0$ di
        $\ell^1$ dan andaikan, demi kontradiksi, bahwa
        $\norm{x^{(k)}}_1 \geq \delta > 0$ sepanjang sebuah
        subbarisan. Tunjukkan lebih dahulu bahwa $x^{(k)}_n \to 0$ untuk
        setiap $n$ (lewat fungsional yang mana?), lalu bangunlah
        secara rekursif indeks $k_1 < k_2 < \cdots$ dan bilangan bulat
        $0 = N_0 < N_1 < N_2 < \cdots$ sedemikian sehingga massa
        $x^{(k_j)}$ memusat pada blok $B_j =
        \intoc{N_{j-1}}{N_j}$:
        \[
        \sum_{n \in B_j}\bigl|x^{(k_j)}_n\bigr| \geq
        \norm{x^{(k_j)}}_1 - \frac\delta{10} .
        \]
  \item (Schur, langkah 2) Definisikan $y \in \ell^\infty$ dengan $y_n =
        \operatorname{sign}\bigl(x^{(k_j)}_n\bigr)$ untuk $n \in
        B_j$. Tunjukkan $\langle x^{(k_j)}, y\rangle \geq
        \norm{x^{(k_j)}}_1 - \frac{2\delta}{10} \geq
        \frac{8\delta}{10}$ lalu turunkan kontradiksi dengan
        $x^{(k)} \rightharpoonup 0$. Simpulkan \emph{teorema
        Schur}: bahwa di $\ell^1$, barisan yang konvergen lemah
        konvergen dalam norma.
  \item Simpulkan bahwa $(e_k)$ tak punya subbarisan yang konvergen
        lemah di $\ell^1$ (sebab satu-satunya calon limitnya adalah
        $0$ secara koordinat --- lalu terapkan Schur), sehingga memperoleh kembali
        kegagalan kekompakan lemah pada pertanyaan 13; lalu selesaikanlah
        paradoks yang tampak itu: di $\ell^1$ kekonvergenan lemah dan norma bagi
        \emph{barisan} berimpit, namun
        \emph{\omterm{def:b3:topology:topology}{topologi}} lemah dan normanya berbeda dan himpunan terbatas
        tetap gagal \omterm{def:b3:topology:compact}{kompak} barisan secara lemah ---
        jadi tak ada kontradiksi, hanya kegagalan kerefleksifan.
  \item (Tabel sintesis) Untuk $E \in \{c_0,\ \ell^1,\ \ell^p\
        (1{<}p{<}\infty),\ \ell^\infty\}$, tabulasikan: \omterm{def:b3:banach:operator}{dualnya};
        keseparabelannya; kerefleksifannya; apakah barisan terbatasnya
        mempunyai subbarisan yang konvergen lemah; dan satu
        sifat khas setiap ruang, yang dibenarkan dalam satu
        baris dari soal ini.
\end{enumerate}

\medskip
\noindent\textbf{Bagian VI --- Pelengkap: titik terdekat,
kekonvergenan rata-rata, dan nilai sebuah limit Banach.}
\begin{enumerate}[resume]
  \item (Titik terdekat: hasil panen kerefleksifan) Misalkan $F$
        sebuah subruang tertutup $\ell^p$ (dengan $1 < p < \infty$) dan
        $x \in \ell^p$. Tunjukkan bahwa $d = \operatorname{dist}(x,
        F)$ \emph{tercapai}: ekstraklah dari sebuah barisan peminimum
        sebuah subbarisan yang konvergen lemah (pertanyaan
        13), tahanlah limit lemahnya tetap di dalam $F$ dengan membangun, lewat
        Hahn--Banach, sebuah fungsional yang lenyap pada $F$ tetapi tidak di
        sebuah titik di luarnya, lalu simpulkan dengan pertanyaan 17.
        Lalu tunjukkan bahwa hak istimewa itu tidak universal: di $c_0$,
        untuk $\Lambda(x) = \sum_n2^{-n}x_n$, buktikan bahwa
        $\norm\Lambda = 1$ tidak tercapai pada bola satuannya,
        tegakkan rumus jaraknya
        $\operatorname{dist}(x, \ker\Lambda) =
        \abs{\Lambda(x)}$, lalu simpulkan bahwa tak ada $x \notin
        \ker\Lambda$ yang mempunyai titik terdekat di hiperbidang
        tertutup $\ker\Lambda$.
  \item (Banach--Saks di $\ell^2$) Misalkan $x^{(k)}
        \rightharpoonup 0$ di $\ell^2_{\R}$ dengan
        $\norm{x^{(k)}}_2 \leq C$. Bangunlah sebuah subbarisan
        $(y_j)$ dengan $\abs{\langle y_i, y_j\rangle} \leq
        \frac1j$ untuk setiap $i < j$, lalu simpulkan
        \[
        \Bigl\lVert\frac{y_1 + \dots + y_m}m\Bigr\rVert_2^2
        \leq \frac{C^2 + 2}m \longrightarrow 0 :
        \]
        jadi setelah ekstraksi, rata-rata Cesàronya konvergen dalam
        \emph{norma}. Periksalah pada $(e_k)$, yang rata-ratanya bernorma
        $\frac1{\sqrt m}$: jadi kekonvergenan lemah, yang tak berguna bagi
        barisannya sendiri (pertanyaan 16), menjadi kekonvergenan norma
        bagi rata-ratanya.
  \item (Nilai sebuah limit Banach) Misalkan $L$ sembarang limit
        Banach (pertanyaan 10) dan $A_mx = \frac1m(x + Sx + \dots
        + S^{m-1}x)$. Tunjukkan $L(A_mx) = L(x)$ dan $\liminf A_mx
        \leq L(x) \leq \limsup A_mx$; lalu simpulkan bahwa \emph{semua}
        limit Banach bersesuaian pada barisan periodik, dengan nilai
        berupa rata-rata atas satu periodenya --- yakni $\frac13$ pada $(1, 0, 0, 1,
        0, 0, \dots)$, sejalan dengan $\frac12$ pada
        pertanyaan 12. Lalu tunjukkan bahwa kesesuaiannya gagal secara umum: untuk
        barisan blok $x$ yang bernilai $1$ pada
        $\intoc{3^{j-1}}{3^j}$ untuk $j$ genap dan $0$ selainnya,
        tunjukkan bahwa rata-rata Cesàronya berayun antara $\leq
        \frac13$ dan $\geq \frac23$, lalu bangunlah dua limit Banach
        $L_\pm$ dengan $L_-(x) \leq \frac13 < \frac23
        \leq L_+(x)$ \emph{(perluaslah dari $c \oplus \R x$ dengan
        nilai ekstrem $\pm$ yang diperbolehkan: periksalah bahwa
        $\Lambda(y + tx) = \lim y + t\,p(x)$ didominasi oleh $p$ yang
        sublinear pada pertanyaan 10)}.

\end{enumerate}
\end{problem}
